\begin{graphillustration}{Пример изменений плоской укладки: удаление немостового $ac$ объединяет две грани; удаление моста $bv$ сохраняет грани и оставляет изолированную вершину $v$.}{fig:13-faces}
\begin{tikzpicture}[font=\small]
\begin{scope}[shift={(0.0,0)}]
\node[gi vertex] (a) at (0,0) {$a$};
\node[gi vertex] (b) at (2,0) {$b$};
\node[gi vertex] (c) at (2,2) {$c$};
\node[gi vertex] (d) at (0,2) {$d$};
\node[gi vertex] (v) at (3,0) {$v$};
\draw[gi edge] (a) -- (b);
\draw[gi edge] (b) -- (c);
\draw[gi edge] (c) -- (d);
\draw[gi edge] (d) -- (a);
\draw[gi edge] (a) -- (c);
\draw[gi edge] (b) -- (v);
\node[gi note] at (1,-0.8) {$G$: $f=3$, $c=1$};
\node[gi note] at (-0.45,1) {$F_0$};
\node[gi note] at (0.5,1.3) {$F_1$};
\node[gi note] at (1.5,0.6) {$F_2$};
\end{scope}
\begin{scope}[shift={(4.3,0)}]
\node[gi vertex] (a) at (0,0) {$a$};
\node[gi vertex] (b) at (2,0) {$b$};
\node[gi vertex] (c) at (2,2) {$c$};
\node[gi vertex] (d) at (0,2) {$d$};
\node[gi vertex] (v) at (3,0) {$v$};
\draw[gi edge] (a) -- (b);
\draw[gi edge] (b) -- (c);
\draw[gi edge] (c) -- (d);
\draw[gi edge] (d) -- (a);
\draw[gi edge] (b) -- (v);
\draw[gi removed] (a) -- (c);
\node[gi note] at (1,-0.8) {$G-ac$: $f=2$, $c=1$};
\node[gi note] at (-0.45,1) {$F_0$};
\node[gi note] at (0.75,1.3) {$F_1$};
\end{scope}
\begin{scope}[shift={(8.6,0)}]
\node[gi vertex] (a) at (0,0) {$a$};
\node[gi vertex] (b) at (2,0) {$b$};
\node[gi vertex] (c) at (2,2) {$c$};
\node[gi vertex] (d) at (0,2) {$d$};
\node[gi vertex] (v) at (3,0) {$v$};
\draw[gi edge] (a) -- (b);
\draw[gi edge] (b) -- (c);
\draw[gi edge] (c) -- (d);
\draw[gi edge] (d) -- (a);
\draw[gi edge] (a) -- (c);
\draw[gi removed] (b) -- (v);
\node[gi note] at (1,-0.8) {$G-bv$: $f=3$, $c=2$};
\node[gi note] at (-0.45,1) {$F_0$};
\node[gi note] at (0.5,1.3) {$F_1$};
\node[gi note] at (1.5,0.6) {$F_2$};
\end{scope}
\end{tikzpicture}
\end{graphillustration}
